\begin{code}[hide]
module slides.popl.sections.evaluation where

open import Data.Product using (_×_ ; _,_)
import Relation.Binary.PropositionalEquality as Eq
open import Word.Base using (Word ; WRel ; [_]ʷ ; ε ; _•_ ; _ᵗ)

infix 4 _~_ _===_

data Gen : Set where
  X-gen S-gen ω-gen : Gen
data C : Set where
  X-cr ε-cr : C
postulate
  Sω : Set
  h  : C → Gen → Word Sω × C
  _~_ : Word Sω × C → Word Sω × C → Set
  by-equal-nf : ∀ {A : Set} → A
data _===_ : WRel Gen where
  order-ω order-S order-X order-SX : ε === ε
  comm : ∀ {gen : Gen} → ε === ε
\end{code}
%% Proof by evaluation (paper §4.4, §4.7, §5.4).

\begin{frame}{Obligations close by evaluation}
\begin{columns}[c]
\begin{column}{0.64\textwidth}
\small A real level of the single-qubit Clifford+$T$ tower: adjoin $X$ to
$\langle\omega\rangle\times\langle S\rangle$, two cosets $\varepsilon$, $X$.
Hypothesis (2), verbatim:
\renewcommand{\AgdaCodeStyle}{\scriptsize}
\begin{code}
h-wd-ax : ∀ c {u t} → u === t → (h ᵗ) c u ~ (h ᵗ) c t
h-wd-ax X-cr {u} {t} order-ω = (by-equal-nf Eq.refl) , Eq.refl
h-wd-ax X-cr {u} {t} order-S = (by-equal-nf Eq.refl) , Eq.refl
h-wd-ax X-cr {u} {t} order-X = (by-equal-nf Eq.refl) , Eq.refl
h-wd-ax X-cr {u} {t} order-SX = (by-equal-nf Eq.refl) , Eq.refl
h-wd-ax X-cr {u} {t} (comm {X-gen}) = (by-equal-nf Eq.refl) , Eq.refl
\end{code}
{\scriptsize \ldots\ nine more; from \texttt{Examples/Amalgamations/CliffordT1.agda}}
\vspace{4pt}

Run the table on both sides of the axiom: the cosets coincide, and the
residues have the \alert{same normal form one level down}.
\end{column}
\begin{column}{0.33\textwidth}
\centering
\begin{tikzpicture}[node distance=4mm]
  \node[bx=prim,text width=38mm,font=\small] (a) {\textbf{coverage checker}\\{\scriptsize no case is missed}};
  \node[bx=ok,text width=38mm,font=\small,below=of a] (b) {\textbf{evaluator}\\{\scriptsize closes every case that is computation}};
  \node[bx=hot,text width=38mm,font=\small,below=of b] (c) {\textbf{you}\\{\scriptsize the few real derivations, which read like the paper's}};
\end{tikzpicture}

\vspace{8pt}
\begin{tabular}{@{}rl@{}}
\textbf{59} & obligations: qubit Clifford+$T$\\
\textbf{72} & one qutrit level\\
\end{tabular}

{\small every one \aid{refl}}
\end{column}
\end{columns}
\end{frame}

\begin{frame}{\textbf{4}\; Indices are chosen for the unifier}
\begin{columns}[c]
\begin{column}{0.48\textwidth}
\aid{\_+\_} recurses on its \emph{first} argument:
\[ 2 + n \;\leadsto\; \mathsf{suc}\,(\mathsf{suc}\,n) \qquad n + 2 \;\;\text{is stuck.} \]
So the indices are \alert{left-biased}: \aid{gate₂} lands in
\aid{Gen (₂₊ n)}, and the $k$-fold shift in \aid{Gen (k + n)}.
\vspace{6pt}

Every index equation the coverage checker meets is then solved by
\alert{injectivity of \aid{suc}}, never by arithmetic --- which is why coset
tables are written by plain pattern matching.
\end{column}
\begin{column}{0.48\textwidth}
\centering
\begin{tikzpicture}[node distance=6mm]
  \node[bx=ok,text width=52mm,align=left,font=\small] (a) {\textbf{left-biased} \hfill\checkmark\\[3pt]
    \aid{Gen (₂₊ n) ≟ Gen (₁₊ m)}\\
    $\leadsto$ $m := \aid{₁₊ n}$, by unification};
  \node[bx=hot,text width=52mm,align=left,font=\small,below=of a] (b) {\textbf{right-biased} \hfill$\times$\\[3pt]
    \aid{Gen (n + 2) ≟ Gen (₁₊ m)}\\
    $\leadsto$ $n + 2 \doteq \mathsf{suc}\;m$: stuck,\\
    a Diophantine side condition};
\end{tikzpicture}
\end{column}
\end{columns}
\end{frame}
